\section{Review Method}
\label{sec:methodology}

This survey covers localization and mapping, traversability and motion control, semantic navigation, multi-robot coordination, and fault management for wheeled mobile robots operating under degraded sensing, terrain, maps, communication, or computation. It includes journal articles and main-conference papers from 2021 to 2026 that propose, evaluate, or benchmark a mechanism for these functions. Studies on road vehicles, other ground robots, or non-ground carriers are included when the inputs of their mechanism exist on a wheeled robot and its assumptions do not rest on the other embodiment. Fig.~\ref{fig:matrix} and Table~\ref{tab:evaluation} count these platform groups separately.

The \CorpusCount{} publications were selected in purposive stages. A language-model-assisted title and abstract screening of robotics journals formed the initial sample, which theme-balanced and targeted searches extended together with works that an internal review of an earlier draft named as missing (Section~\ref{S-supp:selection} of the supplementary material). The corpus is dominated by IEEE publications (\IEEEShare\%) and by works from 2025--2026 (\RecentShare\%). Its counts in Section~\ref{sec:evaluation} therefore describe the selected literature rather than the field. Foundational methods from before 2021 are cited as background outside the coded corpus.

Each work is coded with a published codebook (Section~\ref{S-supp:coding} of the supplementary material). Degradation evidence falls into three classes: a condition injected or varied against a comparator (D1), a condition deliberately selected for its difficulty without manipulation (D2), and an incidental condition that the study reports (D3). A named degradation condition means D1 or D2. A combined attribute such as a task outcome under such a condition counts only when both parts come from the same experiment. A separate language-model session coded a stratified sample of \SampleCount{} works without access to the primary codes, with a median Cohen's kappa of \KappaMedian{} over the double-coded attributes (Section~\ref{S-supp:second-coding} of the supplementary material). This value measures agreement between two language-model coding passes and is not human inter-rater reliability. Reported metrics are interpreted within the platform and protocol of each study, without pooling effect sizes across datasets.
